% ch-a4-19.tex —— 第 4 篇（Emerton--Gee--Hellmann）续片，底本 PDF 396--407 = 印刷页 383--394
% 本片承接上一片末尾（附录 C.1 中段），止于命题 D.1.3 证明中段；续片自 PDF 408 页起。

\ysrc{383}

……$\Hom_{\OO_X}(\OO_X,-)$，二者合起来诱导 $A$-模范畴与拟凝聚 $A$-模范畴之间的等价。这是 Morita 理论所提供的更一般结构的一个实例，下面我们更形式地讨论它。

为进一步阐发这个例子，我们指出：可以把任何交换环 $A$ 视为一个「抽象的」交换算子环，把 $A$-模 $M$ 视为 $A$ 作为算子环的一个具体实现（即通过它在 $M$ 上的作用）。这时注意，$X=\Spec A$ 可以看作在参数化 $A$ 的元素所能取的一切可能的联立特征值系统（因为 $\Spec A$ 的点通过取核，对应于态射 $A\to F$（$F$ 为域）的等价类；两个态射 $A\to E$ 与 $A\to F$ 等价，是指存在 $E$ 与 $F$ 的公共扩域 $K$，使得诱导的态射 $A\to K$ 相合）；因此 $X$ 被称为 $A$ 的「谱」。粗略地说，与 $M$ 相伴的拟凝聚层 $\MM$ 的支集，由那些在 $M$ 中出现的联立特征值系统组成。（更精确地说，$\MM$ 在一点处的纤维，是 $M$ 关于相应特征值系统的余特征空间。）这样，局部化就给出了谱理论的代数解释。

\begin{exam}{C.1.2}\label{thm:a4-19-1}
若取 $A=k[x]$（$k$ 为某个域）并应用上述讨论，我们就重新得到通常经由分析 $k[x]$-模的结构而对单个算子的谱理论所作的推导。
\end{exam}

我们感兴趣的是其他不那么明显的局部化例子。但在每种情形中，正如在局部化 $A$-模这一典范例子中一样，我们希望与模相伴的层的结构（它的纤维、支集等等）能阐明并分离出内部结构中那些若仅用模来考虑则较为晦涩的部分。

\subsection*{C.2.\quad Morita 理论}

Morita 理论——至少我们在这里所意指的——是一个略带笼统意味的名称，指描述各种环上模范畴之间函子（尤其是等价）的理论。我们先回忆该理论的经典（即阿贝尔、即模论）情形，然后再转向导出语境（在导出语境中，该理论的想法容易推广，为研究不必是等价的函子提供灵活而方便的框架）。

\subsubsection*{C.2.1.\quad 阿贝尔情形}

经典（即模论）版本 Morita 理论的一种表述方式是：若 $A$ 与 $B$ 是（不必交换的）结合环，相应的左模范畴为 $A\text{-}\Mod$ 与 $B\text{-}\Mod$，则范畴等价 $A\text{-}\Mod\to B\text{-}\Mod$ 全都形如
\[
M\mapsto P\otimes_A M,
\]
其中 $P$ 是有限生成投射 $B$-模，生成 $B\text{-}\Mod$，并带有同构 $A\xrightarrow{\sim}\End_B(P)$。一个拟逆由
\[
N\mapsto \Hom_B(P,N)
\]
给出。（当然，Morita 理论对这一伴随函子同样适用；我们发现 $Q:=P^\vee:=\Hom_B(P,B)$——它是 $(A,B)$-双模——作为 $A$-模必然是投射的，而拟逆可以改写为 $N\mapsto N\otimes_B Q$。）

\ysrc{384}

更一般地，任何与余极限交换的函子 $F:A\text{-}\Mod\to B\text{-}\Mod$ 都形如
\begin{equation}\tag{C.2.2}\label{eq:a4-19-1}
F:M\mapsto P\otimes_A M
\end{equation}
其中 $P$ 是某个 $(B,A)$-双模。（该双模 $P$ 由在 $A$ 上求值给出。然后对 $M$ 的某个表现求值 $F$，便得到公式 (C.2.2)。）$F$ 的右伴随（由于 $F$ 与余极限交换而存在）显然由 $\Hom_B(P,-)$ 给出。

若 $F$ 进一步全忠实，则我们看到 $A\to\End_B(P)$ 必是同构。若它实际上是等价，则 $P$ 必然继承 $A$ 在 $A\text{-}\Mod$ 中所具有的一切纯范畴性质（紧性、投射性、作为生成元），这就给出经典 Morita 理论的陈述。

\begin{exam}{C.2.3}\label{thm:a4-19-2}
经典 Morita 理论的一个标准例子由 $B=M_n(A)$、$P=A^{\oplus n}$（视为长度为 $n$ 的列向量）的情形给出，带有显然的 $(M_n(A),A)$-双模结构。由于 $B\xrightarrow{\sim}P^{\oplus n}$ 作为其自身上的模，我们看到 $P$ 有限生成且投射，于是 $M\mapsto P\otimes_A M=M^{\oplus n}$ 给出 $A\text{-}\Mod$ 与 $M_n(A)\text{-}\Mod$ 两范畴之间的等价。
\end{exam}

反过来，若令 $Q=A^{\oplus n}$（现在视为长度为 $n$ 的行向量），带有显然的 $(A,M_n(A))$-双模结构，则 $N\mapsto Q\otimes_{M_n(A)}N$ 给出 $M_n(A)\text{-}\Mod$ 与 $A\text{-}\Mod$ 之间的拟逆等价。

\begin{exam}{C.2.4}\label{thm:a4-19-3}
当 $A$ 交换时，可以把上例变成一个局部化的例子：写 $X=\Spec A$，并考虑函子
\[
N\mapsto\OO_X^{\oplus n}\otimes_{M_n(A)}N,
\]
它诱导 $M_n(A)\text{-}\Mod$ 与 $X$ 上拟凝聚层范畴之间的范畴等价。
\end{exam}

当然，我们感兴趣的是研究比交换环上的矩阵环更复杂的非交换环上的模。在这种语境中，一般没有理由指望得到与某个交换环上的模范畴、或与概形上拟凝聚层范畴的范畴等价。正如我们将看到的，更合理的期望是全忠实嵌入；同时，扩展承载层的「空间」$X$ 的可能性（例如允许它是代数叠）也是有意义的。

此外，转向导出语境也更自然。的确，若考虑判定形如 (C.2.2) 的函子何时全忠实（但不必是等价）的问题，就会立即遇到关于 $F$ 的同调性质的问题，而这些在导出世界中考虑更为自然。下面我们将简要讨论导出 Morita 理论，着重于它在研究不必是等价的函子方面所提供的额外灵活性（与经典阿贝尔情形相比）。

\subsubsection*{C.2.5.\quad 稳定 $\infty$-范畴的情形}

设 $A$ 与 $B$ 是 $E_1$-环（$E_1$-环的术语见附录 A.2）。则依 [Lur17, Prop.~7.1.2.4]，连续函子 $\mathrm{LMod}_A\to\mathrm{LMod}_B$ 由与 $(A,B)$-双模作张量给出。特别地，若 $A$ 与 $B$ 是静态的（即通常意义下的结合环），则稳定 $\infty$-范畴的等价 $D(A\text{-}\Mod)\to D(B\text{-}\Mod)$ 全都形如
\begin{equation}\tag{C.2.6}\label{eq:a4-19-2}
M\mapsto T\otimes_A^{\mathbf L}M,
\end{equation}

\ysrc{385}

其中 $T$ 是 $B$-模的完美复形，生成 $D(B\text{-}\Mod)$，并带有 $E_1$-环的同构 $A^{\mathrm{op}}\xrightarrow{\sim}\mathrm{REND}_B(T)$；这里以 $\mathrm{REND}_B(T)$ 记 $E_1$-环 $\Hom_{D(B\text{-}\Mod)}(T,T)$。

若放弃「完美复形 $T$ 生成 $D(B\text{-}\Mod)$」这一假设，但仍假定 $A^{\mathrm{op}}\xrightarrow{\sim}\mathrm{REND}_B(T)$，则 (C.2.6) 仍定义一个连续全忠实函子 $F:D(A\text{-}\Mod)\to D(B\text{-}\Mod)$，其右伴随 $G:D(B\text{-}\Mod)\to D(A\text{-}\Mod)$ 由
\begin{equation}\tag{C.2.7}\label{eq:a4-19-3}
N\mapsto\mathrm{RHom}_B(T,N)
\end{equation}
给出。

更一般地，若 $T$ 是 $B$-模的完美复形、$S$ 是 $A$-模的完美复形，$S$ 生成 $D(A\text{-}\Mod)$，并带有同构 $\mathrm{REND}_A(S)\xrightarrow{\sim}\mathrm{REND}_B(T)=:\mathcal H^{\mathrm{op}}$，则把由 $S$ 给出的 $D(A\text{-}\Mod)$ 与 $\mathcal H\text{-}\Mod$ 的等价、与由 $T$ 给出的从 $\mathcal H\text{-}\Mod$ 到 $D(B\text{-}\Mod)$ 的连续全忠实函子结合起来，我们得到连续全忠实函子 $D(A\text{-}\Mod)\to D(B\text{-}\Mod)$：
\[
M\mapsto T\otimes_{\mathcal H}^{\mathbf L}\mathrm{RHom}_A(S,M),
\]
其右伴随 $D(B\text{-}\Mod)\to D(A\text{-}\Mod)$ 由
\[
N\mapsto S\otimes_{\mathcal H}^{\mathbf L}\mathrm{RHom}_B(T,N)
\]
给出。

当然，这个框架有许多推广。特别地，在正文中我们考虑这样的例子：把 $D(B\text{-}\Mod)$ 换成叠上某个适当的层导出范畴，如下面的例子所示（另见注 6.1.24）。

\begin{rem}{C.2.8}\label{thm:a4-19-4}
用 $\mathcal H$ 记导出自同态环，意在提示（导出的）Hecke 代数。例如见 [Hel23, (1.1)]，那里给出了一个从表示范畴到以这种方式定义的层范畴的函子。
\end{rem}

\begin{rem}{C.2.9}\label{thm:a4-19-5}
在文献中，导出 Hecke 代数常被看作 dga，而非 $E_1$-环。这两种概念之间的对应如下：依 [Lur17, Prop.~7.1.4.6]，若 $A$ 是（静态）交换环，则 $A$ 上 $E_1$-代数的 $\infty$-范畴可以等同于微分分次 $A$-代数的 $\infty$-范畴。例如，若取 $A=\F_p$，Schneider 的导出 Hecke 代数 [Sch15a] 同样既可以看作 $\F_p$-dga，也可以看作 $\F_p$-$E_1$-代数。
\end{rem}

\subsection*{C.3.\quad 非交换现象}

有一些现象出现在非交换环上的模论中，却不出现在交换情形中；在尝试于非交换设定下作局部化时，必须把它们考虑在内。我们简要回顾其中一些，并解释为什么允许在代数叠（而不只是概形）上作局部化有助于缓解这种局面。

非交换环 $A$ 上的模有一个交换语境中不会出现的基本现象：若 $S_1$ 与 $S_2$ 是不同构的单 $A$-模，则尽管 $\Hom_A(S_1,S_2)$ 显然为零，高阶 $\Ext$ 群，即 $\Ext^i_A(S_1,S_2)$（$i>0$），却未必为零。这意味着：如果把 $A\text{-}\Mod$ 以既保持 $\Hom$ 空间、又在高阶 $\Ext$ 群上诱导同构的方式函子式地映入某个层范畴，那么 $S_1$ 与 $S_2$ 的像不能具有不交的支集。

\ysrc{386}

记住，在 Noether 概形\footnote{我们作此假设只是为了保证讨论不至于陷入技术性复杂情形。}上，拟凝聚层范畴中的单对象恰是闭点处的摩天层，因而不同构的单对象具有不交的支集；于是我们看到，例如若 $A$ 是非交换环，容许具有该性质（即 $\Hom_A(S_1,S_2)=0$ 但 $\mathrm{RHom}_A(S_1,S_2)\neq0$）的单模 $S_1$ 与 $S_2$，则 $A\text{-}\Mod$ 范畴当然不可能等价于 Noether 概形上的拟凝聚层范畴。

\begin{rem}{C.3.1}\label{thm:a4-19-6}
上述讨论有助于解释：在 Beilinson--Bernstein 局部化理论中（该理论确实给出某些李代数表示范畴与旗簇上代数 $\mathcal D$-模范畴之间的等价），出现的为什么是拟凝聚 $\mathcal D$-模，而不仅仅是拟凝聚 $\OO$-模。关键在于 Riemann--Hilbert 对应表明 $\mathcal D$-模的行为像反常层：在后者中（作为初步近似）单对象由它们的支集描述（支集是环境概形 $X$ 的某个分层中的一个层），并且支撑在不同分层上的层之间有有趣的扩张类。（最基本的例子是正合列
\[
0\to j_!\mathbf C\to\mathbf C\to i_*\mathbf C\to0,
\]
其中 $j:\mathbb A^1\setminus\{0\}\to\mathbb A^1$、$i:\{0\}\to\mathbb A^1$ 分别是自然的开、闭浸入。）
\end{rem}

若我们只想要全忠实函子，各种（不）可能性就不那么清楚了：此时单 $A$-模不必映到单拟凝聚层（也可能只是：其像的真拟凝聚子层本身不在函子的像中）。若转到导出语境，情形更为模糊，因为此时「单模」与「子模」的概念都没有定义。尽管如此，非交换情形与交换情形在行为上的差异造成了一种张力，使得难以对非交换环上的模构造有趣的局部化函子。缓解这种张力的一种办法，是允许目标范畴为某个代数叠 $\XX$ 上的拟凝聚层范畴，下面我们就来解释。

若 $x\in|\XX|$ 是代数叠 $\XX$（的底拓扑空间）的闭点，则它通常并不对应某个域 $k$ 的闭浸入 $\Spec k\into\XX$，而是（在温和的有限性假设下，例如拟紧对角；见 [Sta, Tag 06RD]，以及 [EG23, App.~E] 的讨论）对应闭浸入 $\mathcal Z\into\XX$，其中 $\mathcal Z$ 是单点带叠。例如，若我们的叠 $\XX$ 在代数闭域 $k$ 上有限型，则对某个 $k$ 上的代数群 $G$ 有 $\mathcal Z=[\Spec k/G]$\tn{底本此处原文作「for some algebraic group over $G$」，语序有误，应为「对某个 $k$ 上的代数群 $G$」；译文据文意译出（原书 p.386）。}。

那么，为固定想法，假设有闭嵌入 $[\Spec k/G]\into\XX$，其中 $k$ 为某个域、$G$ 为 $k$ 上某个代数群，它给出闭点 $x\in|\XX|$。$G$ 的表示对应 $\XX$ 上支撑在 $x$ 的拟凝聚层；特别地，不可约表示对应 $\XX$ 上拟凝聚层范畴中的单对象。在例子中，我们常常可以通过考虑 $\XX$ 中 $x$ 的高阶无穷小邻域来产生这些层的非分裂扩张。

\ysrc{387}

\begin{exam}{C.3.2}\label{thm:a4-19-7}
考虑 $\XX=[\mathbb A^1/\GG_m]$，其中 $\mathbb A^1=\Spec k[x]$、$\GG_m=\Spec k[t^{\pm1}]$，且 $\GG_m$ 通过 $t\cdot x=tx$ 作用在 $\mathbb A^1$ 上。令 $\mathcal I\subset\OO_\XX$ 记切出闭点 $\{x=0\}$ 的理想层。则我们有凝聚层的非分裂短正合列
\[
0\to k(-1)\cong\mathcal I/\mathcal I^2\to\OO_\XX/\mathcal I^2\to\OO_\XX/\mathcal I\cong k(0)\to0
\]
（这里照常把 $\GG_m$-表示与分次 $k$-模等同起来）。
\end{exam}

简言之，在代数叠上的凝聚层范畴中，可以找到容许非分裂扩张的不同构单对象。这使得代数叠上的凝聚层范畴成为定义在非交换环上模范畴上的局部化函子的合适目标。

\subsection*{C.4.\quad 例子}

下面给出 Morita 理论与局部化的一些例子，从很简单的例子逐步构建到出现在 $\mathrm{GL}_2(\Q_p)$ 表示论中的例子。全篇都在某个交换环 $k$ 上工作。

\begin{exam}{C.4.1}\label{thm:a4-19-8}
如我们在上面例 C.1.2 中所见，范畴 $k[x]\text{-}\Mod$ 等价于 $\Qcoh(\mathbb A^1)$；这个 Morita 等价由结构层 $\OO_{\mathbb A^1}$ 中介（它当然是其整体截面环，即 $k[x]$ 上的模）。稍微不那么平凡地，作为例 C.2.3 的特例，这些范畴也等价于矩阵代数
\[
\begin{pmatrix}k[x]&k[x]\\k[x]&k[x]\end{pmatrix}
\]
的左模范畴，经由双模 $\OO_{\mathbb A^1}\oplus\OO_{\mathbb A^1}$。
\end{exam}

\begin{exam}{C.4.2}\label{thm:a4-19-9}
改而考虑代数
\[
A:=\begin{pmatrix}k[x]&k[x]\\xk[x]&k[x]\end{pmatrix}.
\]
它的左模范畴不与任何交换环上的模范畴 Morita 等价。然而，存在从 $D(A\text{-}\Mod)$ 到 $D(\Qcoh(\XX))$ 的全忠实函子，其中
\[
\XX:=[(\Spec k[x,y]/(xy))/\GG_m],
\]
$\GG_m$-作用由\footnote{此处作用公式中出现 $t^2$ 而非 $t$，对本例没有本质重要性；它唯一的意义在于给出一个与范畴 Langlands 直接相联系的例子，如下文所注。}
\[
t\cdot(x,y)=(x,t^2y)
\]
给出，而函子由与
\[
P:=\OO_\XX\oplus(\OO_\XX/y)
\]
作（导出！）张量给出。

为使这一点有意义，我们需要 $P$ 是右 $A$-模。我们断言事实上
\[
\mathrm{REND}_{\OO_\XX}(P)=\End_{\OO_\XX}(P)=A^{\mathrm{op}},
\]
于是函子 $P\otimes_A-$ 全忠实。

一般地，对于阿贝尔范畴中的对象 $X\oplus Y$，有
{\small
\[
\mathrm{REND}(X\oplus Y)=\begin{pmatrix}\mathrm{REND}(X)&\mathrm{RHom}(Y,X)\\\mathrm{RHom}(Y,X)&\mathrm{REND}(Y)\end{pmatrix},
\]
}
\tn{底本此矩阵左下角印作 $\mathrm{RHom}(Y,X)$；按下一行 $\mathrm{REND}(X\oplus Y)^{\mathrm{op}}$ 的矩阵（右上角为 $\mathrm{RHom}(X,Y)$），此处应为 $\mathrm{RHom}(X,Y)$。译文照录（原书 p.387）。}

\ysrc{388}

它通过矩阵从左侧作用在 $X\oplus Y$ 上（此时把 $X\oplus Y$ 看作列向量）。于是
{\small
\[
\mathrm{REND}(X\oplus Y)^{\mathrm{op}}=\begin{pmatrix}\mathrm{REND}(X)&\mathrm{RHom}(X,Y)\\\mathrm{RHom}(Y,X)&\mathrm{REND}(Y)\end{pmatrix},
\]
}
它通过矩阵从右侧作用在 $X\oplus Y$ 上（此时把 $X\oplus Y$ 看作行向量）。

在我们的例子中，$\XX$ 是一个仿射概形模线性约化群所得的商，具体地是模环面 $\GG_m$ 的商，因此 $\XX$ 上的整体截面函子是正合的；等价地，$\OO_\XX$ 投射，且 $\mathrm{RHom}_{\OO_\XX}(\OO_\XX,-)=\Hom_{\OO_\XX}(\OO_\XX,-)$ 就是整体截面函子。特别地
\[
\mathrm{RHom}_{\OO_\XX}(\OO_\XX,\OO_\XX)=\mathrm{RHom}_{\OO_\XX}(\OO_\XX,\OO_\XX/y)=k[x].
\]
为计算 $\mathrm{RHom}_{\OO_\XX}(\OO_\XX/y,-)$，我们使用投射消解
\[
\cdots\xrightarrow{y}\OO_\XX(-2)\xrightarrow{x}\OO_\XX(-2)\xrightarrow{y}\OO_\XX\to\OO_\XX/y\to0.
\]
（这里照常，商叠 $\XX$ 上可逆层的记号对应于 $\Z$-分次模的记号，约定为 $\Gr^iM(n)=\Gr^{i+n}M$。）

我们有：当 $n=0$ 时 $\Hom_{\OO_\XX}(\OO_\XX(-n),\OO_\XX)=k[x]$，当 $n>0$ 时它 $=k\cdot y^n$；类似地，当 $n=0$ 时 $\Hom_{\OO_\XX}(\OO_\XX(-n),\OO_\XX/y)=k[x]$，当 $n>0$ 时它 $=0$。由此推出，$\mathrm{RHom}_{\OO_\XX}(\OO_\XX/y,\OO_\XX)$ 由复形
\[
0\to k[x]\to k\xrightarrow{0}k\xrightarrow{\sim}k\xrightarrow{0}k\to\cdots
\]
计算（其中态射 $k[x]\to k$ 由在 $0$ 处取值给出），而 $\mathrm{RHom}_{\OO_\XX}(\OO_\XX/y,\OO_\XX/y)$ 由复形
\[
0\to k[x]\to0\to0\to\cdots,
\]
计算，从而
\[
\mathrm{RHom}_{\OO_\XX}(\OO_\XX/y,\OO_\XX)=\Hom_{\OO_\XX}(\OO_\XX/y,\OO_\XX)=xk[x],
\]
并且
\[
\mathrm{RHom}_{\OO_\XX}(\OO_\XX/y,\OO_\XX/y)=\Hom_{\OO_\XX}(\OO_\XX/y,\OO_\XX/y)=k[x].
\]
于是确实 $\mathrm{REND}_{\OO_\XX}(P)^{\mathrm{op}}=A$。

再把这个例子继续一点：写 $M=k[x]\oplus xk[x]$（看作列向量，带有显然的 $A$ 左作用）与 $N=k[x]\oplus k[x]$（同样看作列向量，带有显然的 $A$ 左作用），则作为左 $A$-模有 $A=M\oplus N$，于是 $M$ 与 $N$ 在 $A$ 上都投射，且 $P\otimes_AM=\OO_\XX$，而 $P\otimes_AN=\OO_\XX/y$。

显然的包含
\begin{equation}\tag{C.4.3}\label{eq:a4-19-4}
M\into N
\end{equation}
有余核 $S_1:=0\oplus k$（其 $A$-模结构来自把元素看作列向量并让 $x$ 经由 $0$ 作用），而态射
\begin{equation}\tag{C.4.4}\label{eq:a4-19-5}
N\xrightarrow{x}M
\end{equation}
有余核 $S_2:=k\oplus0$（它的元素也看作列向量，其 $A$-模结构同样由 $x$ 经由 $0$ 作用给出）。两个 $A$-模 $S_1,S_2$ 都是单模且不同构（并且在同构意义下，正是 $A$ 的中心 $k[x]$ 中元素 $x$ 经由 $0$ 作用的两个单 $A$-模）。商 $M/xM$ 是 $S_2$ 被 $S_1$ 的非分裂扩张，而 $N/xN$ 是 $S_1$ 被 $S_2$ 的非分裂扩张。

\ysrc{389}

若转到 $\XX$ 上的层，包含 (C.4.3) 对应满射 $\OO_\XX\to\OO_\XX/y$，于是 $S_1$ 对应此态射的锥，即 $(\OO_\XX/x)(-2)$。态射 (C.4.4) 对应态射 $\OO_\XX/y\xrightarrow{x}\OO_\XX$，于是 $S_2$ 对应此态射的锥，即 $\OO_\XX/x$。这样我们看到 §C.3 所讨论的可能性之一例：两个不同构的单模容许它们之间的非分裂扩张，对应着不同构但支集相同的凝聚层。
\end{exam}

\begin{rem}{C.4.5}\label{thm:a4-19-10}
例 C.4.2 与 $\mathrm{PGL}_2(\Q_p)$ 的 $\ell$-进（$\ell\neq p$）表示论密切相关，例如见 [Hel23, Rem.~4.43]。事实上，$\mathrm{PGL}_2(\Q_p)$ 的 Iwahori Hecke 代数 $\mathcal H_I$ 是一个矩阵代数
{\small
\[
\mathcal H_I=\begin{pmatrix}\Q_\ell[T]&\Q_\ell[T]\\(T^2-(p+1)^2)\Q_\ell[T]&\Q_\ell[T]\end{pmatrix},
\]}
且存在全忠实函子（由凝聚 Springer 层决定），从（左）$\mathcal H$-模的导出范畴到模叠 $\XX_2^{\mathrm{unip}}$ 上的 ind-凝聚层范畴，该模叠参数化行列式为分圆特征之逆的二维幂幺分歧 Weil--Deligne 表示。

这个叠有三个不可约成分：非分歧成分，以及两个由二次扭相联系、被称为「Steinberg」的成分。两个 Steinberg 成分互不相交，但每个都与非分歧成分相交，交点分别是表示 $1\oplus\varepsilon^{-1}$ 及其二次扭。在这两个交点中任一个的邻域内，我们可以把 $\XX_2^{\mathrm{unip}}$ 与（上述叠 $\XX$ 的奇点处的对应邻域）等同起来。此时轨迹 $y=0$ 对应 Steinberg 成分，而 $P$ 对应凝聚 Springer 层（例如见 [Zhu25a, Ex.~4.4.4]）。

模 $M$ 对应 Steinberg 表示（或其二次扭），其相伴层为 $\OO_\XX/y$，即 Steinberg 成分的结构层；而 $N$ 对应平凡表示（或其二次扭），其相伴层为 $(\OO_\XX/y)(-2)[1]$，即 Steinberg 成分结构层的一个扭的平移。
\end{rem}

\begin{rem}{C.4.6}\label{thm:a4-19-11}
注意例 C.4.2 中的对象 $P$ 不是完美复形，但它是凝聚的；因此已经在这个例子里我们就看到，在 Langlands 参数的叠上考虑 ind-凝聚层（而不是拟凝聚层）是自然的。
\end{rem}

下面这个例 C.4.2 的变体在 $\mathrm{GL}_2(\Q_p)$ 的 $p$-进 Langlands 对应中起类似作用，§7.3 对此有描述。

\begin{exam}{C.4.7}\label{thm:a4-19-12}
设 $A$ 是（交换）环，考虑分次环
\[
B:=A[x,y],
\]
其中分次结构由分别给 $x$ 与 $y$ 权 $2$ 与 $-2$ 定义。以 $\mathcal B$ 记分次 $B$-模的范畴（等价地，商叠 $[\Spec B/\GG_m]$ 上的拟凝聚层范畴）。若 $M$ 是分次模、$n\in\Z$，则照常以 $M(n)$ 记这样的分次模：其底模与 $M$ 相同，且 $\Gr^iM(n)=\Gr^{i+n}M$。
\end{exam}

\ysrc{390}

我们考虑模 $P:=B(1)\oplus B(-1)$，并定义 $R:=\End_{\mathcal B}(P)^{\mathrm{op}}$；于是 $R$ 从右侧作用在 $P$ 上。具体计算给出
{\small
\[
\begin{aligned}
R&=\begin{pmatrix}\End(B(1))&\Hom(B(1),B(-1))\\\Hom(B(-1),B(1))&\End(B(-1))\end{pmatrix}
=\begin{pmatrix}\Gr^0B&\Gr^0B(-2)\\\Gr^0B(2)&\Gr^0B\end{pmatrix}\\
&=\begin{pmatrix}\Gr^0B&\Gr^{-2}B\\\Gr^2B&\Gr^0B\end{pmatrix}
=\begin{pmatrix}A[xy]&yA[xy]\\xA[xy]&A[xy]\end{pmatrix},
\end{aligned}
\]}
它通过通常的矩阵从右侧乘以行向量作用在 $P$ 上。

当然，$P$ 是 $\mathcal B$ 的投射对象。下面这个引理则不那么显然。

\begin{lem}{C.4.8}\label{thm:a4-19-13}
$P$ 作为右 $R$-模是投射的。
\end{lem}

\begin{proof}
我们首先注意，作为自身上的右模，有 $R=R_1\oplus R_2$，其中 $R_1=A[xy]\oplus yA[xy]$（把 $R$ 如上视为矩阵序（matrix order）时的第一行），$R_2=xA[xy]\oplus A[xy]$（$R$ 的第二行）。$R_1$ 与 $R_2$ 各自都是投射的（作为 $R$ 的直和项）。

其次注意 $P=\bigoplus_{m=-\infty}^{\infty}\Gr^{1+2m}P$；由于 $R$ 作为 $P$ 的分次自同态环作用，这是 $R$-模的同构。若 $1+2m\ge1$（即 $m\ge0$），则 $\Gr^{1+2m}P=x^mA[xy]\oplus x^{m+1}A[xy]\xrightarrow{\sim}A[xy]\oplus xA[xy]=R_2$；而若 $1+2m\le-1$（即 $m\le-1$），则 $\Gr^{1+2m}P=y^{-m-1}A[xy]\oplus y^{-m}A[xy]\xrightarrow{\sim}A[xy]\oplus yA[xy]=R_1$。于是 $P$ 是投射 $R$-模的直和，从而本身是投射 $R$-模。
\end{proof}

我们有如下直接推论。

\begin{prop}{C.4.9}\label{thm:a4-19-14}
函子 $N\mapsto P\otimes_R N$ 给出从 $R$-模的阿贝尔范畴到 $\mathcal B$ 的全忠实正合嵌入，其本质像等于 $\mathcal B$ 中由 $P$ 生成的全子范畴 $\mathcal C$。
\end{prop}

事实上，由于 $P$ 在 $B$ 上投射，故 $\mathrm{RHom}_B(P,P)=R$，因此上述命题有导出类比（依照 §C.2.5 的讨论）。

\begin{prop}{C.4.10}\label{thm:a4-19-15}
函子 $N\mapsto P\otimes_R N$ 给出带 $t$-结构的稳定 $\infty$-范畴的全忠实 $t$-正合嵌入 $D(R)\to D(\mathcal B)$，其本质像等于 $D(\mathcal B)$ 中由 $P$ 生成的全子范畴。
\end{prop}

\subsubsection*{C.4.11.\quad 刻画一个像}

正如已经指出的，例 C.4.7 在 $\mathrm{GL}_2(\Q_p)$ 情形的范畴 $p$-进 Langlands 对应中有一个应用（见 §7.5）。事实上，在这个应用中我们将研究此函子的一个变体：在其中我们使用稳定 $\infty$-范畴，并额外施加一个支集条件。

现在精确地描述我们的设定。首先，我们保持例 C.4.7 的记号，并进一步假设环 $A$ 是 Noether 环。我们还引入一些额外记号：以 $\mathcal R$ 记 $R$-模的阿贝尔范畴，以 $\mathcal R_0$ 记 $\mathcal R$ 中由这样的模组成的全子范畴：其中每个元素都被 $xy$ 的某个幂零化。类似地，以 $\mathcal B_0$ 记分次 $B$-模范畴 $\mathcal B$ 中满足同一条件（即每个元素都被 $xy$ 的某个幂零化）的模组成的子范畴。

命题 B.1.17（以及它对分次环的显然变体）表明，包含 $\mathcal R_0\into R\text{-}\Mod$ 与 $\mathcal B_0\into\mathcal B$ 诱导出等价
\begin{equation}\tag{C.4.12}\label{eq:a4-19-6}
D^+(\mathcal R_0)\xrightarrow{\sim}D^+_{(xy)}(R)
\end{equation}

\ysrc{391}

以及
\begin{equation}\tag{C.4.13}\label{eq:a4-19-7}
D^+(\mathcal B_0)\xrightarrow{\sim}D^+_{(xy)}(\mathcal B).
\end{equation}
（这里的下标 $(xy)$ 表示由这样的对象组成的全子范畴：其各阶上同调都被 $xy$ 的某个幂逐元素零化。）

\begin{lem}{C.4.14}\label{thm:a4-19-16}
等价 (C.4.12) 与 (C.4.13) 延拓为等价 $D(\mathcal R_0)\to D_{(xy)}(R)$ 与 $D(\mathcal B_0)\to D_{(xy)}(\mathcal B)$。
\end{lem}

\begin{proof}
这将由命题 A.7.3(2) 推出，只需证明：包含 $\mathcal R_0\into R\text{-}\Mod$ 与 $\mathcal B_0\into\mathcal B$ 的伴随函子具有有界上同调维数。然而无论哪种情形，该伴随函子的导出函子都由 $X\mapsto\operatorname*{colim}_n\bigl(X\xrightarrow{(xy)^n}X\bigr)$ 计算，因此其上同调幅度为 $1$。
\end{proof}

命题 C.4.10 的 $t$-正合全忠实函子 $N\mapsto P\otimes_R N$ 显然限制为 $t$-正合全忠实函子 $D_{(xy)}(R)\to D_{(xy)}(\mathcal B)$\tn{底本此处原文作「a $t$-exact functor fully faithful functor」，重出一词，疑为笔误（原书 p.391）。}；由引理 C.4.14，我们也可以把它视为 $t$-正合全忠实函子 $D(\mathcal R_0)\to D(\mathcal B_0)$。我们将用某个半正交分解来刻画这个函子的像，下面就来解释。

伴随函子 $\Hom_{\mathcal B}(P,-):D_{(xy)}(\mathcal B)\to D_{(xy)}(R)$ 有更显式的描述（记住 $P:=B(1)\oplus B(-1)$）：
\[
M\mapsto\Gr^{-1}M\oplus\Gr^1M.
\]
（右端带有显然的 $R$-结构：每个直和项自然地是 $A[xy]$-模；乘以 $x$ 把第一个直和项送到第二个；乘以 $y$ 把第二个直和项送到第一个。）

\begin{prop}{C.4.15}\label{thm:a4-19-17}
以 $\mathcal K$ 记 $D_{(xy)}(\mathcal B)$ 中由满足 $\Gr^{-1}M=\Gr^1M=0$ 的对象 $M$ 组成的全子范畴。则 $P\otimes_R-$ 的像等于 ${}^{\perp}\mathcal K$。
\end{prop}

\begin{proof}
这来自引理 A.8.4(6)。
\end{proof}

\subsubsection*{C.4.16.\quad $\mathcal K$ 的一组生成元}

为了范畴 $p$-进 Langlands 的应用，我们希望用一组生成元来给出核 $\mathcal K$ 的另一种描述。为此，我们稍微改变语境：依照 §A.6 所述的程序，把无界稳定 $\infty$-范畴 $D(\mathcal R_0)$ 与 $D(\mathcal B_0)$ 换成它们的正则化 $\IndCoh(\mathcal R_0)$ 与 $\IndCoh(\mathcal B_0)$。

函子 $P\otimes_R-$ 及其伴随 $\Hom_{\mathcal B}(P,-)$ 都正合且连续（后者连续是因为 $P$ 作为 $B$-模有限生成且投射），并且保持紧对象。因此它们在 $\IndCoh(\mathcal R_0)$ 与 $\IndCoh(\mathcal B_0)$ 之间诱导一对伴随函子；上文在 $D(\mathcal R_0)$、$D(\mathcal B_0)$ 语境中证明的结果，在 ind-凝聚语境中也成立。

范畴 $\mathcal B_0$ 分裂为乘积 $\mathcal B_0^+\times\mathcal B_0^-$，其中 $\mathcal B_0^+$（相应地 $\mathcal B_0^-$）由分次片仅支撑在偶（相应地奇）次数的对象组成。相应地，$\IndCoh(\mathcal B_0)$ 分解为乘积 $\IndCoh(\mathcal B_0^+)\times\IndCoh(\mathcal B_0^-)$。

函子 $P\otimes_R-$ 显然经过 $\IndCoh(\mathcal B_0^-)$ 分解；我们以 $\mathcal D_0^-$ 记它在后一范畴中的本质像。我们还考虑类似的函子 $B\otimes_{A[xy]}-:\IndCoh(\mathcal A_0)\to\IndCoh(\mathcal B_0^+)$，其中 $\mathcal A_0$ 记 $A[xy]\text{-}\Mod$ 中由被 $xy$ 的某个幂逐元素零化的模组成的子范畴。我们以 $\mathcal D_0^+$ 记后一函子在 $\IndCoh(\mathcal B_0^+)$ 中的本质像。

\ysrc{392}

我们经由同构 $B/(y)\xrightarrow{\sim}A[x]$ 把 $A[x]$ 视为分次 $B$-模，对 $A[y]$ 类似。同样，我们经由同构 $B_i/(x,y)\xrightarrow{\sim}A$ 把 $A$ 视为分次 $B$-模\tn{底本此处印作 $B_i/(x,y)$，按上下文应为 $B/(x,y)$（下标 $i$ 疑为误植）；译文照录（原书 p.392）。}。于是 $A[x]$、$A[y]$、$A$ 三者都是 $\mathcal B_0$ 的对象。

使用分次消解
\begin{equation}\tag{C.4.17}\label{eq:a4-19-8}
0\to B(2)\xrightarrow{y}B\to A[x]\to0,
\end{equation}
我们看到（对任意选定的扭 $n$），$\mathrm{RHom}_{\mathcal B}(A[x],A[x](n))$ 由上同调次数 $0$ 与 $1$ 中的复形
\[
\Gr^0(A[x](n))\xrightarrow{0}\Gr^0(A[x](n-2))
\]
计算，而 $\mathrm{RHom}_{\mathcal B}(A[x],A[y](n))$ 由上同调次数 $0$ 与 $1$ 中的复形
\[
0\to\Gr^0A(n-2)
\]
计算。

类似地，我们得到 $\mathrm{RHom}_{\mathcal B}(A[y],A[y](n))$ 由上同调次数 $0$ 与 $1$ 中的复形
\[
\Gr^0(A[y](n))\xrightarrow{0}\Gr^0(A[y](n+2))
\]
计算，而 $\mathrm{RHom}_{\mathcal B}(A[y],A[x](n))$ 由上同调次数 $0$ 与 $1$ 中的复形
\[
0\to\Gr^0A(n+2)
\]
计算。

于是我们看到，对象族 $\{A[x](-2),A[x](-4),\dots\}$ 与 $\{A[y](2),A[y](4),\dots\}$ 各自都在 $\IndCoh(\mathcal B_0^+)$ 中构成定义 A.8.10 意义下的弱例外集，并且彼此正交。由引理 A.8.11，它们诱导出 $\IndCoh(\mathcal B_0^+)$ 中由它们生成的全子范畴 $\mathcal E_0^+$ 的半正交分解。

类似地，对象族 $\{A[x](-3),A[x](-5),\dots\}$ 与 $\{A[y](3),A[y](5),\dots\}$ 各自都在 $\IndCoh(\mathcal B_0^-)$ 中构成弱例外集，并且彼此正交。它们同样诱导出 $\IndCoh(\mathcal B_0^-)$ 中由它们生成的全子范畴 $\mathcal E_0^-$ 的半正交分解。

\begin{prop}{C.4.18}\label{thm:a4-19-18}
\begin{enumerate}[label=(\arabic*)]
\item $\mathcal D_0^+$ 是 $\IndCoh(\mathcal B_0^+)$ 中由对象 $B/(xy)$ 生成的余完备全子范畴。
\item $\mathcal E_0^+=(\mathcal D_0^+)^{\perp}$。
\item $\mathcal D_0^-$ 是 $\IndCoh(\mathcal B_0^+)$ 中由对象 $B/(xy)(1)$ 与 $B/(xy)(-1)$ 生成的余完备全子范畴\tn{底本此处印作 $\IndCoh(\mathcal B_0^+)$；按上文定义，$\mathcal D_0^-$ 是 $\IndCoh(\mathcal B_0^-)$ 中的本质像，故此处应为 $\IndCoh(\mathcal B_0^-)$。译文照录（原书 p.392）。}。
\item $\mathcal E_0^-=(\mathcal D_0^-)^{\perp}$。
\end{enumerate}
\end{prop}

\begin{proof}
余完备稳定 $\infty$-范畴 $\IndCoh(\mathcal A_0)$ 由 $A:=A[xy]/(xy)$ 生成。于是 $\mathcal D_0^+$（它被定义为该范畴在函子 $B\otimes_{A[xy]}-$ 下的像）由 $A$ 的像，即 $B/(xy)$ 生成。这证明了 (1)。

引理 A.8.4 与命题 C.4.15 的 ind-凝聚版本表明，$(\mathcal D_0^+)^{\perp}$ 由满足 $\Gr^0M=0$ 的对象 $M$ 组成。另一方面，引理 A.8.6 表明，$\mathcal E_0^+$ 可以描述为 $D(\mathcal B_0^+)$ 中包含对象集 $X$ 的最小余完备稳定子范畴。由于 $X$ 中每个对象都满足 $\Gr M_0=0$\tn{底本此处（以及下一段）印作 $\Gr M_0$，且下一段又印作 $(\mathcal D_0^+)$；据上下文，两处应分别为 $\Gr^0M$ 与 $(\mathcal D_0^+)^{\perp}$。译文照录（原书 p.392）。}，我们于是得到 $\mathcal E_0^+\subseteq(\mathcal D_0^+)^{\perp}$。

现在，考虑 (C.4.17) 的 $-2$ 扭，并利用 $(\mathcal D_0^+)$ 中对象满足 $\Gr M_0=0$ 这一事实，可以看出：对 $(\mathcal D_0^+)^{\perp}$ 中的对象 $M$ 有 $\mathrm{RHom}_{\mathcal B}(A[x](-2),M)=\Gr^2M$。特别地，$\mathrm{RHom}_{\mathcal B}(A[x](-2),A[x](-2))=A$。

\ysrc{393}

于是 $N\mapsto A[x](-2)\otimes_AN$ 是全忠实函子 $\IndCoh(A)\to(\mathcal D_0^+)^{\perp}$，其伴随由 $M\mapsto\Gr^2M$ 给出。特别地，$A[x](-2)\otimes_A\Gr^2M\to M$ 的锥的 $\Gr^2$ 为零。归纳地继续，依次考虑各个 $A[x](-n)$，然后还考虑 $A[y](n)$（$n$ 取正偶数值），我们发现：若 $M$ 是 $(\mathcal E_0^+)^{\perp}\cap(\mathcal D_0^+)^{\perp}$ 的对象，则对一切 $i$ 都有 $\Gr^iM=0$，从而 $M=0$。因此 $\mathcal E_0^+=(\mathcal D_0^+)^{\perp}$，这证明了 (2)。

(3) 与 (4) 的证明与 (1)、(2) 完全类似，故从略。
\end{proof}

\section*{附录 D.\quad 拓扑阿贝尔群、环与模}

\subsection*{D.1.\quad 泛函分析}

设 $\{G_n\}$ 是一列线性拓扑化\footnote{即 $0$ 有一组由子群组成的邻域基。}的 Hausdorff 拓扑阿贝尔群，带有闭嵌入 $G_n\into G_{n+1}$。于是可以定义归纳极限 $G:=\varinjlim_n G_n$，并赋予归纳极限线性拓扑（即这样的拓扑：子集为开，当且仅当它在每个 $G_n$ 中的原像为开）；这仍是一个线性拓扑化拓扑阿贝尔群。我们回忆 $G$ 满足的一些基本性质；它们类比于各种语境（LF 空间、紧型空间等）中一列局部凸拓扑空间的归纳极限的标准性质，而我们的论证是对 [RR73, Ch.~VII] 中论证的改编。

\begin{rem}{D.1.1}\label{thm:a4-19-19}
\begin{enumerate}[label=(\arabic*)]
\item 为简化记号，我们经由典范嵌入把每个 $G_n$ 看作 $G$ 的子集。下文将会看到，这个群的嵌入还是拓扑嵌入，因此这不应引起混淆。
\item 按定义，$G$ 的子群 $U$ 为开，是指对每个 $n$，拉回 $G_n\cap U$ 都是 $G_n$ 的开子群。构造这类开子群的一种典型方式是：对每个 $n$ 给定 $G_n$ 的开子群 $V_n$，并令 $U_n:=\sum_{i=1}^n V_i$。则 $U_n$ 仍是 $G_n$ 的开子群，且显然 $U_n\subseteq U_{n+1}$，于是 $U:=\bigcup_n U_n=\sum_{i=1}^\infty V_i$ 是 $G$ 的子群。由于 $U_n\subseteq G_n\cap U$，可见 $U$ 是 $G$ 的开子群。若 $V_n$ 还满足额外条件 $G_n\cap V_{n+1}\subseteq V_n$，则可以看出 $U_n=G_n\cap U_{n+1}$，从而 $U_n=G_n\cap U$。
\end{enumerate}

若给定开子群 $V_n$，则可以归纳地构造开子群 $V'_n\subseteq V_n$，使得进一步有 $G_n\cap V'_{n+1}\subseteq V'_n$。具体地说，假设对 $i=1,\dots,n$ 已构造出 $V'_i$；由于 $G_n\into G_{n+1}$ 是拓扑嵌入，我们可以找到开子群 $W_{n+1}$，使得 $G_n\cap W_{n+1}\subseteq V'_n$。然后令 $V'_{n+1}:=V_{n+1}\cap W_{n+1}$。
\end{rem}

\begin{lem}{D.1.2}\label{thm:a4-19-20}
若 $X\subseteq G$ 具有性质：对每个 $n$，$X\cap G_n$ 有限，则作为 $G$ 的子集，$X$ 既闭又离散。
\end{lem}

\begin{proof}
我们先证明 $X$ 离散。为此取 $x\in X$；我们需要找到 $x$ 在 $G$ 中的一个邻域，它与 $X\setminus\{x\}$ 不相交。为简化记号，把 $X$ 换成 $X-x$，把 $x$ 换成 $0$。由于对每个 $n$，$G_n\cap(X\setminus\{0\})$ 有限，我们可以归纳地选取 $G_n$ 的开子群 $V_n$，使得 $U_n:=V_1+\cdots+V_n$ 与 $G_n\cap(X\setminus\{0\})$ 不相交。于是（依照注 D.1.1(2)）$U:=\bigcup U_n$ 是 $G$ 的开子群，且与 $X\setminus\{0\}$ 不相交。

\ysrc{394}

现在假设 $x$ 是 $G\setminus X$ 的元素。则 $X':=X\cup\{x\}$ 仍满足引理的假设，因此由已证部分，它是离散的。于是可以找到 $x$ 的一个邻域与 $X$ 不相交，所以 $x$ 不在 $X$ 的闭包中。这就证明了 $X$ 确实闭。
\end{proof}

\begin{prop}{D.1.3}\label{thm:a4-19-21}
\begin{enumerate}[label=(\arabic*)]
\item 典范映射 $G_n\into G$ 中的每一个都是闭嵌入。
\item $G$ 是 Hausdorff 的；若每个 $G_n$ 都完备，则 $G$ 完备。
\item $G$ 的任意紧子集都包含在某个 $G_n$ 中。
\item 若 $H$ 是 $G$ 的子群，并令 $H_n:=G_n\cap H$，则自然映射 $\varinjlim_n H_n\to H$ 与 $\varinjlim_n G_n/H_n\to G/H$ 各自都是拓扑同构。此外，$H$ 在 $G$ 中闭，当且仅当对每个 $n$，$H_n$ 在 $G_n$ 中闭。
\end{enumerate}
\end{prop}

\begin{proof}
固定 $n_0$，令 $U_{n_0}$ 是 $G_{n_0}$ 的任意开子群。若对 $n\le n_0$ 令 $V_n:=U_{n_0}\cap G_n$，并对 $n>n_0$ 选取 $V_n$ 使得 $V_n\cap G_{n-1}\subseteq V_{n-1}$，则我们可以依照注 D.1.1(2) 找到开子群 $U\subset G$，使得 $U\cap G_{n_0}=U_{n_0}$。于是 $G_{n_0}$ 的拓扑由 $G$ 的拓扑诱导，这证明了 (1)。

为看出 $G$ 是 Hausdorff 的，只需证明：若 $g\in G\setminus\{0\}$，则可以找到不含 $g$ 的 $G$ 的开子群。选取 $n_0$ 使得 $g\in G_{n_0}\setminus\{0\}$，并令 $U_{n_0}$ 是 $G_{n_0}$ 的任意开子群，使得 $g\notin U_{n_0}$。由上一段的结果，可以找到 $G$ 的开子群 $U$ 使得 $G_{n_0}\cap U=U_{n_0}$。于是 $g\notin U$，正如所需。这证明了 (2) 的第一个断言。

现在假设每个 $G_n$ 都完备，并以 $\widehat G$ 记 $G$ 的完备化。恒等态射 $G\to\widehat G$ 是拓扑嵌入（因为如上所见 $G$ 是 Hausdorff 的）。由于每个 $G_n$ 完备，复合嵌入 $G_n\into G\into\widehat G$ 把 $G_n$ 实现为 $\widehat G$ 的闭子群。若 $G$ 不完备，则可以找到 $x\in\widehat G\setminus G$；又因每个 $G_n$ 在 $\widehat G$ 中闭（如上所注）且不含 $x$，可以找到 $\widehat G$ 的开子群 $W_n$，使得 $x+W_n$ 与 $G_n$ 不相交；我们还可以选取 $W_n$ 使得 $W_{n+1}\subseteq W_n$。

现在令 $U:=\sum_{n=1}^\infty(W_n\cap G_n)$；由注 D.1.1，这是 $G$ 的开子群。由完备化的构造，它的完备化 $\widehat U$（即它在 $\widehat G$ 中的闭包）是 $\widehat G$ 的开子群。由于 $G$ 在 $\widehat G$ 中稠密，我们看到 $(x+\widehat U)\cap G$ 非空，从而对某个 $n$，$(x+\widehat U)\cap G_n$ 非空。另一方面，有
{\small
\[
\begin{aligned}
\widehat U&\subseteq U+W_n=\sum_{i=1}^{\infty}(W_i\cap G_i)+W_n\\
&=\sum_{i\le n}(W_i\cap G_i)+\sum_{i>n}(W_i\cap G_i)+W_n
\subseteq\sum_{i\le n}G_n+\sum_{i>n}W_i+W_n\\
&=G_n+W_{n+1}+W_n=G_n+W_n,
\end{aligned}
\]}
于是 $(x+G_n+W_n)\cap G_n$ 非空，从而 $(x+W_n)\cap G_n$ 非空，这与我们对 $W_n$ 的选择矛盾。这表明必有 $G=\widehat G$，从而完成了 (2) 的证明。

令 $X\subseteq G$ 是紧子集；由于 $G$ 是 Hausdorff 的（由已证部分），$X$ 是闭的。假设对每个 $n$ 都有 $X\not\subseteq G_n$\tn{底本此处印作 $X\subsetneq G_n$，按上下文应为 $X\not\subseteq G_n$（否则不能取到 $x_n\in X$ 使 $x_n\notin G_n$）。译文据文意译出（原书 p.394）。}；则可以找到 $X$ 中元素组成的序列 $x_n$，使得 $x_n\notin G_n$。于是 $\bigcup_{n=1}^\infty\{x_n\}$ 是一个无限子集。

% 本片止于印刷页 394 页末（命题 D.1.3 证明中段）；proof 环境在下一片闭合。
% 续片 ch-a4-20.tex 自印 395 起接排，并在 (4) 证完后 \end{proof} 排 □。
